% ===================================================================
%  TÁBLA Uimh. 12 — An Stáisiún. — Na hIarnróid. — Na Longa.
%  Traduction 2026 d'apres texto/io, controlee sur texto/fr, texto/en
%  et texto/es. Consigne : voir texto/ga/10-tabla-01.tex.
%
%  MEME OBSERVATION QU'AU TABLEAU 10, ET LE PARTAGE EST AUSSI NET :
%  « traein », « vaigín », « inneall », « tolgán », « teileagraf »,
%  « seamafór », « tender » sont venus avec le rail, tout faits, au
%  XIXe siecle. Ce que l'irlandais fournit en propre, c'est ce que le
%  rail a trouve en place et n'a pas remplace : « fead » le sifflet,
%  « roth » la roue, « doras » la portiere, « mála » le sac, « cófra »
%  le coffre, « ciseán » le panier — mots de la maison et de la main,
%  appris dedans et sans maitre.
%
%  ZERO INVERSION ATTENDUE, ET ZERO RELEVEE. Ce tableau enumere plus
%  qu'il ne construit : la plupart de ses renvois se suivent dans une
%  liste, et une liste ne s'inverse pas. Les trois passages qui
%  auraient pu couter — « la kaldiero (56) di la lokomotivo », « la
%  fairisto (50), sur sua tendro (49) », « remorkeri (82) potenta
%  tiris pezoza kargo-bateli (83) » — se rangent d'eux-memes : le
%  genitif irlandais postpose, et le VSO met ici le porteur de renvoi
%  en sujet, devant un objet qui en porte un lui aussi mais qui vient
%  DEJA apres dans l'ido.
% ===================================================================

%%K t12-apar-1 apar
\VUsaut{4.0mm}
\VUtitre{15.0pt}{60}{TÁBLA Uimh. 12}
\VUsaut{2.4mm}
\VUpk{11.4pt}{40}{An Stáisiún. --- Na hIarnróid. --- Na Longa.}
\VUsaut{3.4mm}
\textsc{Nóta.} --- \textit{Is éagsúil na hamchláir a bhíonn i ngach stáisiún, agus
mar sin glacaimid mar shampla amanna} \VUgras{an tábla Fhrancaigh}.

%%K t12-01-1 p
1. --- Rinne mé turas tríd an nGearmáin le déanaí. Roimh imeacht dom cheannaigh mé
\VUgras{clár ama} \textsuperscript{(15)} chun uaireanta imeachta na dtraenacha a
fháil amach. Nuair a chinntigh mé gurbh í an traein ba áisiúla an ceann a d'imigh ó
Pháras ag a naoi a chlog tráthnóna agus a shroicheann Strasbourg ag a haon a chlog
trí nóiméad déag, chuir mé mo thuras in ord, agus lá arna mhárach thug mé orm féin
dul chun an stáisiúin.

%%K t12-02-1 p
2. --- Nuair a chuaigh mé isteach sa halla imeachta mór, taobh thiar de
\VUgras{shagart} tuaithe aosta \textsuperscript{(2)}, a bhí ag iompar a \VUgras{mhála
taistil} \textsuperscript{(3)} ina láimh, bhí a lán \VUgras{taistealaithe}
\textsuperscript{(1)} tagtha cheana.

%%K t12-02-2 p
Ó theastaigh uaim \VUgras{teachtaireacht (teileagram)} \textsuperscript{(5)} a chur,
chuir mé \VUgras{cigire} \textsuperscript{(6)} ag taispeáint \VUgras{oifig an
teileagraif} \textsuperscript{(4)} dom.

%%K t12-03-1 p
3. --- Sula ndeachaigh mé isteach sa \VUgras{seomra feithimh} \textsuperscript{(7)}
chuaigh mé ag fáil \VUgras{ticéid} \textsuperscript{(9)} fillte dom féin.

%%K t12-04-1 p
4. --- Níor fhan mé i bhfad ag an \VUgras{bhfuinneoigín} \textsuperscript{(8)}, mar
ní raibh mórán daoine ann. Bhí \VUgras{saighdiúir} \textsuperscript{(10)}, a bhí ag
imeacht ar saoire, sásta a áit a thabhairt dom, ionas go raibh mo dhóthain ama agam
chun roinnt eolais a iarraidh ar \VUgras{mháistir an stáisiúin} \textsuperscript{(13)},
a bhí ag comhrá leis an \VUgras{gcoimisinéir} \textsuperscript{(12)} faireachais agus
leis an \VUgras{ngarda} \textsuperscript{(11)} a bhí ar dualgas.

%%K t12-tit-1 sub
\VUpk{10.2pt}{40}{Clárú an bhagáiste.}
\VUsaut{3.0mm}

%%K t12-05-1 p
5. --- Tar éis nuachtán a cheannach ag an \VUgras{leabharlann}
\textsuperscript{(14)}, chuir mé mo chuid \VUgras{bagáiste} \textsuperscript{(17)} á
mheá, bagáiste a thug \VUgras{póirtéir} \textsuperscript{(16)} isteach ar
\VUgras{thralaí pacáistí} \textsuperscript{(19)}; chuir \VUgras{an cléireach}
\textsuperscript{(22)} a bhí i mbun na \VUgras{meá} \textsuperscript{(21)} in iúl dom
go raibh cúig chileagram is tríocha i mo chófra; b'in \VUgras{rómheáchan} de chúig
chileagram, agus chaithfinn a íoc dá bharr ag \VUgras{fuinneoigín} an bhagáiste
\textsuperscript{(23)}.

%%K t12-06-1 p
6. --- Ó bhí mé róluath, bhí uain agam chun breathnú ar shiúl na dtaistealaithe sa
stáisiún. Chuir cuid acu a mbagáiste beag i dtaisce nó thóg ar ais é ag an
\VUgras{taisceadán bagáiste} \textsuperscript{(25)}; chuaigh cuid eile i gcomhairle
leis na \VUgras{cláir ama} \textsuperscript{(26)}. Bhí na taistealaithe a tháinig
isteach ag brostú chun an \VUgras{doras amach} \textsuperscript{(32)} agus a
\VUgras{mála} \textsuperscript{(24)} ina láimh acu. D'fhan roinnt acu ag an
\VUgras{oifig bhagáiste} \textsuperscript{(27)} agus thaispeáin siad a \VUgras{ticéad}
\textsuperscript{(29)} chun a gcófraí, a \VUgras{gcistí} \textsuperscript{(28)} nó a
\VUgras{gciseáin} \textsuperscript{(30)} a fháil ar ais.

%%K t12-07-1 p
7. --- Scrúdaigh \VUgras{oifigigh an mháil} \textsuperscript{(33)} na pacáistí go
haireach féachaint an raibh aon rud le dearbhú.

%%K t12-08-1 p
8. --- Chuaigh cuid de na taistealaithe chun an \VUgras{bhuféidh}
\textsuperscript{(34)}, áit a raibh \VUgras{freastalaí} \textsuperscript{(35)} go
dícheallach á bhfreastal. Caithfidh siad brostú, mar ní fhanfaidh an \VUgras{traein}
\textsuperscript{(36)}, ar luastraein í, i bhfad sa stáisiún.

%%K t12-09-1 p
9. --- Ansin chuaigh mé chun ardáin an imeachta agus chuardaigh mé \VUgras{carráiste}
díreach \textsuperscript{(37)} le nach gcuirfí orm carráiste a athrú i lár an turais.
Chuaigh mé suas i gcarráiste pasáiste ina bhfuil \VUgras{urrann}
\textsuperscript{(38)} do lucht tobac agus urrann curtha ar leataobh do « mhná
uaisle gan tionlacan ».

%%K t12-tit-2 sub
\VUpk{10.2pt}{40}{Imeacht san oíche.}
\VUsaut{3.0mm}

%%K t12-10-1 p
10. --- Tar éis dul suas ar an \VUgras{gcéim} \textsuperscript{(40)}, an
\VUgras{doras} \textsuperscript{(39)} a dhúnadh, an \VUgras{dallóg}
\textsuperscript{(41)} a chornadh agus mo bhagáiste beag a chur ar an
\VUgras{líonra} \textsuperscript{(43)}, shuigh mé ar an \VUgras{mbinse}
\textsuperscript{(42)}. Nuair a chonaic mé an \VUgras{lampadóir}
\textsuperscript{(44)} ag lasadh na lampaí, smaoinigh mé go gcaithfinn oíche a
chaitheamh sa charráiste, agus ghlaoigh mé ar an oifigeach a bhí i mbun ligean na
\VUgras{gceannadhart} \textsuperscript{(45)} chun ceann acu a iarraidh air.

%%K t12-11-1 p
11. --- Tháinig stiúrthóir na traenach ag seiceáil na dticéad. D'fhiafraigh mé de cén
fáth nach raibh imithe againn cheana, mar ba é an t-am ceart é. D'fhreagair sé go
raibh \VUgras{ceannaire na traenach} \textsuperscript{(46)} gafa le rothair a chur sa
\VUgras{veain} \textsuperscript{(47)}. Ina theannta sin ní raibh na téitheoirí coise
\textsuperscript{(48)} go léir athraithe fós.

%%K t12-12-1 p
12. --- Ach ní fada go n-imeoimis, mar bhí \VUgras{an tiománaí teallaigh}
\textsuperscript{(50)}, ar a \VUgras{thender} \textsuperscript{(49)}, ag tógáil
guail chun tine \VUgras{an innill} \textsuperscript{(54)} a ghríosadh, agus ní raibh
\VUgras{an t-innealtóir} \textsuperscript{(51)} ag fanacht ach le comhartha
\VUgras{leas-mháistir an stáisiúin} \textsuperscript{(52)}, a bhí ar an \VUgras{ardán}
\textsuperscript{(53)}.

%%K t12-13-1 p
13. --- Bhí \VUgras{coire} \textsuperscript{(56)} an innill ag drantán go bodhar;
bhí \VUgras{an simléar} \textsuperscript{(55)} ag caitheamh amach tuilte deataigh
measctha le \VUgras{gal} \textsuperscript{(60)}. Bhúir \VUgras{fead} ghéar
\textsuperscript{(59)} agus thosaigh na \VUgras{loiní} \textsuperscript{(58)} ag
bogadh, ag tiomáint \VUgras{rothaí} \textsuperscript{(57)} an innill throm.

%%K t12-tit-3 sub
\VUsaut{3.0mm}
\VUpk{10.2pt}{40}{Ar Siúl!}
\VUsaut{3.0mm}

%%K t12-14-1 p
14. --- Ghéaraigh luas na traenach agus í ag rith ar na \VUgras{ráillí}
\textsuperscript{(64)}; d'imigh na \VUgras{hardáin} \textsuperscript{(65)} as radharc;
ghabhamar thar na \VUgras{leithris} \textsuperscript{(66)}, agus thar shraith
carráistí a bhí ar stad ar an \VUgras{líne} eile \textsuperscript{(63)}:
\VUgras{carráistí codlata} \textsuperscript{(67)}, \VUgras{carráiste bia}
\textsuperscript{(68)}, \VUgras{carráiste poist} \textsuperscript{(69)}.

%%K t12-noto-1 noto
\VUnotes{143.2mm}{%
\textit{(*) Nóta.} --- \textit{Is é atá sa chur síos ar an turas, ar ndóigh, an
teorainn mar a bhí sí roimh an gcogadh.}%
}

%%K t12-15-1 p
15. --- Ansin ghabh \VUgras{stáisiún na n-earraí} \textsuperscript{(71)} thar ár
súile. Faoi na sciobóil bhí oifigigh ag luchtú na \VUgras{n-earraí}
\textsuperscript{(72)} a bhí le cur amach ar thraein mhall. Bhí \VUgras{scuad fear}
\textsuperscript{(76)} in aice leis an \VUgras{umar uisce} \textsuperscript{(73)} ag
ainliú \VUgras{vaigíní ainmhithe} \textsuperscript{(75)} agus \VUgras{vaigíní
ardáin} \textsuperscript{(79)} chun \VUgras{traein earraí} \textsuperscript{(80)} a
chur le chéile.

%%K t12-16-1 p
16. --- Ghabhamar thar an \VUgras{ngabhal} deireanach \textsuperscript{(78)}, mar a
raibh fear na dtolgán ina sheasamh; ghabhamar thar an \VUgras{diosca comhartha}
\textsuperscript{(86)} agus d'imigh an stáisiún as radharc i gcéin.

%%K t12-17-1 p
17. --- Níorbh fhada gur thrasnaíomar \VUgras{droichead iarainn}
\textsuperscript{(74)} atá thar an \VUgras{abhainn} \textsuperscript{(81)}. Tar éis an
droichead sin a thrasnú lean muid an abhainn, ar a raibh \VUgras{tugaí}
\textsuperscript{(82)} cumhachtacha ag tarraingt \VUgras{báirsí} troma
\textsuperscript{(83)}. Chonaiceamar fós \VUgras{bád paisinéirí}
\textsuperscript{(84)} agus é ag teacht le taobh \VUgras{ponntúin}
\textsuperscript{(85)}.

%%K t12-18-1 p
18. --- Tar éis gabháil thar roinnt stáisiún ar luas lasrach, chuamar trí roinnt
\VUgras{tolláin} \textsuperscript{(87)} agus d'fhágamar inár ndiaidh \VUgras{tithe
beaga} gan áireamh \textsuperscript{(88)} lucht \VUgras{coimeádta na mbacainní}.
Le luascadh na traenach do mo bhréagadh, thit mé i mo chodladh go trom.

%%K t12-tit-4 sub
\VUsaut{3.0mm}
\VUpk{10.2pt}{40}{Stáisiún Deutsch-Avricourt.}
\VUsaut{3.0mm}

%%K t12-19-1 p
19. --- Dhúisigh guth oifigigh go tobann mé agus é ag glaoch: « Deutsch-Avricourt!
\textsuperscript{(*)}. Amach libh go léir! » Rinneadh cigireacht an chustaim agus
gnásanna leadránacha uile na teorann. B'éigean dom m'uaireadóir póca a chur in
oiriúint don \VUgras{chlog mór} \textsuperscript{(89)} sa stáisiún, a bhí caoga cúig
nóiméad chun tosaigh; ar éigean múscailte dom, is ar éigean a d'aithin mé
\VUgras{an tsnáthaid mhór} \textsuperscript{(90)} thar an \VUgras{gceann beag}
\textsuperscript{(91)}; bhí \VUgras{aghaidh an chloig} \textsuperscript{(92)} féin
trí chéile ar fad de bharr cheo na maidine.

%%K t12-20-1 p
20. --- Fad is a bhí lucht an chustaim i mbun a gcigireachta, léigh mé agus mé ag
méanfach na \VUgras{fógraí} \textsuperscript{(93)} atá ag maisiú bhallaí an
stáisiúin. D'ith mé bricfeasta chun an t-am a chur thart, chuaigh mé i gcomhairle le
léarscáil an \VUgras{líonra} Ghearmánaigh \textsuperscript{(94)} féachaint cén fad a bhí
fós idir mé agus Strasbourg, agus chuaigh mé ar ais i mo charráiste, agus mé
neartaithe ag an dóchas go sroichfinn ceann scríbe mo thurais gan mhoill.
